\textbf{Storms (ERA5 mean sea-level pressure).} We retain ERA5 mean sea-level pressure for 1999-11-17--2000-01-14, sampled every 12 hours over the North Atlantic and Europe (25--80\textdegree N, 40\textdegree W--50\textdegree E)~\cite{Hersbach:2020:ERA5}. A $49\times49$ equal-area grid and 250\,km Gaussian smoothing define join-tree minima. Smoothing is performed on a buffered domain; only extrema of that extended field inside the window are retained, and other cropped basins are filled to their spill levels. These are processed-field minima, not independently validated cyclone centres.

\textbf{Wildfires (fire radiative power).} The Wildfires sequence uses Australian fire observations on 10-km source cells from the November--December 2019 sequence~\cite{Franke:2021:1DP}. Starting on 2019-11-01, 60 daily observations are summed into 30 two-day frames, whose start dates end on 2019-12-29. Active source cells are accumulated on a $64\times64$ layout grid, followed by $\log(1+\mathrm{FRP})$. A weak smooth background removes zero plateaus. The smallest common scale from the declared ladder keeping at most twelve maxima is 2.5 layout cells (50.7169\,km); evaluated frames contain six to eleven leaves. Split-tree maxima describe processed fire-intensity structures, rather than independently tracked individual fires.

\textbf{Radiation fields.} Daily NOAA OLR CDR (PSL Ver01Rev02)~\cite{NOAA:OLR} covers 2024-06-01--2024-07-30 over 15\textdegree S--20\textdegree N, 60--150\textdegree E. Join-tree minima represent low radiation, not validated convective systems. A $49\times49$ equal-area grid uses five-cell buffers, priority-flood cropping, and no added imputation; a global-extremum fallback is logged. The fixed ladder (100, 150, 200, 250, 350, 500, 700, 1000\,km) selects the smallest common scale with at most eight extrema, independently of performance: 500\,km. Source hashes and preprocessing trials are recorded.

\textbf{PM2.5 concentration.} The user-identified integrated China V2 snapshot (10-km grid, DOI 10.57760/sciencedb.35623)~\cite{Dong:2026:PM25} supplies all 366 files from 2000-01-01 through 2000-12-31 over 28--38\textdegree N, 102--114\textdegree E. Files carry 10:30 timestamps without an explicit averaging interval; we do not label them daily averages. The $49\times49$ buffered equal-area grid, priority-flood cropping, and fixed 100\,km smoothing retain the earlier processing settings. The annual input has one to seven split-tree leaves. No frame is removed or selected by performance. Values retain $\mu$g\,m$^{-3}$ units. The automatic direction, tracking, and fixed sequence-wide scalar range are recomputed on the annual population. Metadata record upstream linear-then-nearest filling; we add no imputation. Per-file hashes identify the local snapshot. The earlier 30-day January and four-month comparisons remain exploratory controls, not the current main population.

\textbf{Ring.} We reproduce the archived TMTM author generator and workspace for the Ring control~\cite{Koepp:2023:TMTM,Franke:2021:1DP}: 40 frames on the endpoint-inclusive $14\times14$ grid over $[0,210]^2$, without added smoothing, interpolation or noise. The generator changes a Gaussian radial profile over time. The extracted split tree has one to seven maxima because a sampled ring is not a single isolated extremum. The source commit, generator/workspace hashes and float32 property values are recorded in \path{experiments/ring/dataset.py}. The main RG-MTM comparison uses the same automatic fixed-direction reference as the real sequences; the analytic ring radius is not silently substituted for that reference.

\textbf{Suite boundaries.} The five sequences total 614 frames. The archived Storms 2013/14 sequence (source identifier \texttt{ERA5-2014}) is an independent season, not an alias of the main Storms sequence. The earlier four-month PM2.5 audit is a historical selection control; the primary comparison now covers the full year. The 64-frame perpendicular-bisector experiment is a separate synthetic control; none is added to the main total without an explicit protocol change. All methods share processed fields and overlap-based leaf matching, which isolates layout differences within this extraction but does not validate physical identities.

\textbf{Volumetric and static inputs in the common tables.} Tangaroa uses all 201 original-grid $300\times180\times120$ velocity-magnitude volumes and 10\% persistence, retaining 18--93 join-tree leaves. Cylinder computes three-dimensional $Q$ on $192\times64\times48$ velocities before the $135\times64\times48$ crop, retaining 4--75 split-tree leaves across 102 frames at 2\% per-frame persistence. Nucleon is one unchanged $64\times64$ author slice with 13 split-tree leaves; all R outputs coincide with P. Benzene evaluates a declared OPLS-AA point-charge potential on 21 static $61\times61$ sections, retaining 5--10 join-tree leaves with 0.05\% per-section persistence. Sections S6 and S9 provide the complete input and extraction protocols. These static sections are not time samples.
